L'isòtop de poloni ($Z = 84$) amb un temps de semidesintegració més llarg és el Po-210, que el té de 138 dies. Es desintegra per emissió d'una partícula alfa i origina un isòtop estable de plom (Pb).

\begin{apartat}{1,25}
Escriviu la desintegració del Po-210. Si l'activitat inicial per unitat de massa del Po-210 és d'$1,66\cdot 10^{14}\ \mathrm{Bq/g}$, quina serà l'activitat de 5~mg d'aquest element al cap d'una setmana?
\end{apartat}

\begin{apartat}{1,25}
Els nuclis d'heli que es produeixen en les desintegracions alfa no triguen a captar dos electrons. Suposem que es forma un àtom d'heli amb dos passos ben diferenciats: primerament, es transporta des d'una distància molt gran un primer electró a $0,6\cdot 10^{-10}\ \mathrm{m}$ de la partícula alfa i es manté allà. Posteriorment, un segon electró es porta a l'altra banda de la partícula alfa a $0,6\cdot 10^{-10}\ \mathrm{m}$. La configuració final es mostra a la figura.

\figura[0.5]{fig1.png}

Calculeu el treball efectuat pel camp elèctric en cada un dels dos passos. Quina és l'energia potencial electroestàtica de la configuració final?
\end{apartat}

\dades{%
  $|e| = 1,602\times 10^{-19}\ \mathrm{C}$\\
  $k = \frac{1}{4\pi\varepsilon_0} = 8,99\times 10^{9}\ \mathrm{N\cdot m^2\cdot C^{-2}}$}
